\documentclass[11pt,a4paper]{article}

% ============================================================
%  QIMMA -- COURS : Structure d'espace vectoriel
%  2ème Bac Sciences Mathématiques (A et B)
%  Standard exhaustif : définitions formelles + démonstrations
%  Basé sur templates/qimma-template.tex (charte Qimma)
% ============================================================

\usepackage[french]{babel}
\usepackage[utf8]{inputenc}
\usepackage[T1]{fontenc}
\usepackage{lmodern}
\usepackage[a4paper,margin=2cm,top=2.3cm,bottom=2.6cm]{geometry}
\usepackage{amsmath,amssymb}
\usepackage{xcolor}
\usepackage{graphicx}
\usepackage{tikz}
\usetikzlibrary{shadows,calc}
\usepackage[most]{tcolorbox}
\usepackage{titlesec}
\usepackage{fancyhdr}
\usepackage{enumitem}
\usepackage{pifont}
\usepackage{array}
\usepackage{colortbl}
\usepackage{hyperref}

% ------------------- CHARTE QIMMA -------------------
\definecolor{qteal}{HTML}{0d9488}
\definecolor{qtealdark}{HTML}{134e4a}
\definecolor{qamber}{HTML}{fbbf24}
\definecolor{qambertext}{HTML}{92400e}
\definecolor{ink}{HTML}{1f2937}
\definecolor{inkgray}{HTML}{6b7280}
\definecolor{tealpale}{HTML}{e6f5f3}
\colorlet{colDef}{qteal}
\definecolor{colProp}{HTML}{0f766e}
\definecolor{colEx}{HTML}{b45309}
\definecolor{colRem}{HTML}{9d174d}
\color{ink}

\graphicspath{{../../../../assets/}}
\newcommand{\logofile}{../../../../assets/qimma-logo.pdf}

% ------------------- METADONNEES -------------------
\newcommand{\matiere}{Mathématiques}
\newcommand{\niveau}{2\up{ème} Bac -- Sciences Mathématiques}
\newcommand{\chapitretitre}{Structure d'espace vectoriel}

% ------------------- EN-TETE / PIED DE PAGE -------------------
\pagestyle{fancy}
\fancyhf{}
\renewcommand{\headrulewidth}{0pt}
\renewcommand{\footrulewidth}{0.5pt}
\fancyfoot[L]{\raisebox{-0.15cm}{\includegraphics[height=0.35cm]{\logofile}}~\small\sffamily\color{inkgray}Qimma}
\fancyfoot[C]{\small\sffamily\color{inkgray}\matiere\ -- \niveau}
\fancyfoot[R]{\small\sffamily\color{qtealdark}\thepage}

% ------------------- TITRES DE SECTION -------------------
\titleformat{\section}
  {\normalfont\sffamily\Large\bfseries\color{qtealdark}}
  {\tikz[baseline=-3.2pt]{\node[circle,fill=qteal,text=white,inner sep=0pt,minimum size=8mm]{\sffamily\bfseries\thesection};}}
  {10pt}{}
  [\vspace{-2mm}{\color{qamber}\titlerule[1.6pt]}\vspace{2mm}]
\titlespacing*{\section}{0pt}{16pt}{10pt}

\titleformat{\subsection}
  {\normalfont\sffamily\large\bfseries\color{qtealdark}}
  {\color{qteal}\ding{212}}{6pt}{}
\titlespacing*{\subsection}{0pt}{12pt}{6pt}

% ------------------- BOITES PEDAGOGIQUES -------------------
\newtcolorbox{fiche}[3][]{
  enhanced, breakable,
  colback=white, colframe=#2,
  boxrule=1pt, arc=2mm,
  top=7mm, left=4mm, right=4mm, bottom=3mm,
  drop shadow={black!25, shadow xshift=1mm, shadow yshift=-1mm},
  attach boxed title to top left={xshift=4mm, yshift=-3mm, yshifttext=-1mm},
  boxed title style={colback=#2, arc=1.5mm, boxrule=0pt, left=2mm, right=2mm},
  coltitle=white, fonttitle=\sffamily\bfseries,
  title={#3}, #1
}
\newenvironment{definition}[1][Définition]
  {\begin{fiche}{colDef}{\ding{110}~~#1}}{\end{fiche}}
\newenvironment{propriete}[1][Propriété]
  {\begin{fiche}{colProp}{\ding{72}~~#1}}{\end{fiche}}
\newenvironment{exemple}[1][Exemple]
  {\begin{fiche}{colEx}{\ding{217}~~#1}}{\end{fiche}}
\newenvironment{remarque}[1][Remarque]
  {\begin{fiche}{colRem}{\ding{43}~~#1}}{\end{fiche}}

% Démonstration (hors boîte, style discret)
\newenvironment{demo}
  {\par\smallskip\noindent{\sffamily\bfseries\color{qtealdark}Démonstration.}\ \itshape}
  {\ \normalfont\hfill$\blacksquare$\par\medskip}

\hypersetup{colorlinks=true, linkcolor=qtealdark, urlcolor=qtealdark}

% ============================================================
\begin{document}

% ------------------- BANNIERE DE TITRE -------------------
\begin{tcolorbox}[
  enhanced, boxrule=0pt, arc=4mm,
  interior style={left color=qtealdark, right color=qteal},
  top=8mm, bottom=8mm, left=10mm, right=32mm,
  drop shadow={black!35, shadow xshift=1.5mm, shadow yshift=-1.5mm},
  before skip=0pt, after skip=9mm,
  overlay={
    \node[anchor=north west] at ([xshift=6mm,yshift=-6mm]frame.north west)
      {\includegraphics[height=1.25cm]{\logofile}};
    \node[anchor=north east, fill=qamber, text=qambertext, rounded corners=2pt,
          inner sep=4pt, font=\sffamily\bfseries\small] at ([xshift=-6mm,yshift=-6mm]frame.north east) {COURS};
  }
]
\hspace{1.6cm}\centering
{\sffamily\bfseries\color{white}\fontsize{23}{27}\selectfont \chapitretitre}\\[3mm]
{\sffamily\color{white!90}\large \matiere\ \textbullet\ \niveau}
\end{tcolorbox}

% ------------------- OBJECTIFS -------------------
\begin{tcolorbox}[
  enhanced, colback=tealpale, colframe=qteal, boxrule=0.8pt, arc=2mm,
  left=5mm, right=5mm, top=3mm, bottom=3mm,
  title={\sffamily\bfseries\color{qtealdark}\ding{51}~~Objectifs du chapitre},
  coltitle=qtealdark, colbacktitle=tealpale, boxrule=0.8pt,
  attach title to upper={\par\vspace{1mm}}
]
\begin{itemize}[leftmargin=6mm, label=\ding{212}, itemsep=1mm]
  \item Reconnaître que l'ensemble des vecteurs du plan (puis de l'espace) muni de $+$ et de la multiplication par un scalaire est un espace vectoriel réel.
  \item Manipuler combinaisons linéaires, vecteurs colinéaires, familles libres et liées.
  \item Démontrer qu'une partie est (ou n'est pas) un sous-espace vectoriel à l'aide de la caractérisation pratique.
  \item Reconnaître une base du plan et de l'espace, et exploiter les coordonnées associées.
  \item Faire le lien entre cette structure abstraite et les manipulations vectorielles déjà pratiquées en géométrie.
\end{itemize}
\end{tcolorbox}

\vspace{4mm}

% ------------------- SECTION 1 -------------------
\section{L'ensemble des vecteurs comme espace vectoriel}

\begin{definition}[Espace vectoriel réel]
Un ensemble $E$, muni d'une addition interne $+$ et d'une multiplication externe par les réels $\cdot:\mathbb{R}\times E\to E$, est un \textbf{espace vectoriel réel} (ou $\mathbb{R}$-espace vectoriel) lorsque :
\begin{enumerate}[leftmargin=7mm, itemsep=0.5mm]
  \item $(E,+)$ est un \textbf{groupe commutatif} (vu au chapitre précédent), de neutre noté $\vec0$ ;
  \item pour tous $\lambda,\mu\in\mathbb{R}$, $\vec u,\vec v\in E$ : $\lambda\cdot(\vec u+\vec v)=\lambda\cdot\vec u+\lambda\cdot\vec v$ ; $(\lambda+\mu)\cdot\vec u=\lambda\cdot\vec u+\mu\cdot\vec u$ ; $\lambda\cdot(\mu\cdot\vec u)=(\lambda\mu)\cdot\vec u$ ; $1\cdot\vec u=\vec u$.
\end{enumerate}
Les éléments de $E$ sont appelés \textbf{vecteurs}, ceux de $\mathbb{R}$ des \textbf{scalaires}.
\end{definition}

\begin{propriete}[L'ensemble des vecteurs du plan (puis de l'espace) est un espace vectoriel]
L'ensemble $\mathcal{V}$ des vecteurs du plan, muni de l'addition vectorielle (relation de Chasles) et de la multiplication par un réel, vérifie tous les axiomes ci-dessus : c'est un \textbf{espace vectoriel réel}. Il en va de même pour l'ensemble des vecteurs de l'espace.
\end{propriete}

\begin{remarque}[Ce que change ce point de vue]
Tu manipules déjà des combinaisons de vecteurs depuis la 1\up{ère} Bac (relation de Chasles, $\overrightarrow{AB}+\overrightarrow{BC}=\overrightarrow{AC}$, multiplication par un scalaire). Ce chapitre donne le \textbf{cadre théorique} qui justifie pourquoi ces manipulations (règles de calcul, coordonnées, alignement) fonctionnent toujours de la même façon, quel que soit le contexte géométrique.
\end{remarque}

\begin{propriete}[Règles de calcul dans un espace vectoriel]
Pour tous $\lambda\in\mathbb{R}$, $\vec u\in E$ :
\[
  0\cdot\vec u = \vec0, \qquad \lambda\cdot\vec0=\vec0, \qquad (-1)\cdot\vec u=-\vec u, \qquad \lambda\cdot\vec u=\vec0 \iff \lambda=0 \text{ ou } \vec u=\vec0.
\]
\end{propriete}

\begin{demo}
$0\cdot\vec u=\vec0$ : $0\cdot\vec u=(0+0)\cdot\vec u=0\cdot\vec u+0\cdot\vec u$. En simplifiant (groupe $(E,+)$, on ajoute le symétrique de $0\cdot\vec u$ des deux côtés) : $\vec0=0\cdot\vec u$.
\end{demo}

% ------------------- SECTION 2 -------------------
\section{Combinaisons linéaires, colinéarité}

\begin{definition}[Combinaison linéaire]
Une \textbf{combinaison linéaire} des vecteurs $\vec u_1,\ldots,\vec u_n$ est tout vecteur de la forme
\[
  \lambda_1\vec u_1+\lambda_2\vec u_2+\cdots+\lambda_n\vec u_n, \qquad \lambda_1,\ldots,\lambda_n\in\mathbb{R}.
\]
\end{definition}

\begin{definition}[Vecteurs colinéaires]
Deux vecteurs $\vec u,\vec v$ sont \textbf{colinéaires} lorsque l'un est combinaison linéaire de l'autre : il existe $\lambda\in\mathbb{R}$ tel que $\vec v=\lambda\vec u$ (ou $\vec u=\vec0$).
\end{definition}

\begin{definition}[Famille libre, famille liée]
Une famille $(\vec u_1,\ldots,\vec u_n)$ est \textbf{libre} (ou les vecteurs sont \textbf{linéairement indépendants}) lorsque
\[
  \lambda_1\vec u_1+\cdots+\lambda_n\vec u_n=\vec0 \implies \lambda_1=\cdots=\lambda_n=0.
\]
Une famille qui n'est pas libre est dite \textbf{liée}.
\end{definition}

\begin{propriete}[Caractérisation dans le plan et l'espace]
\begin{itemize}[leftmargin=6mm, itemsep=0.5mm]
  \item Deux vecteurs du plan $(\vec u,\vec v)$ forment une famille \textbf{liée} $\iff$ ils sont colinéaires.
  \item Trois vecteurs de l'espace $(\vec u,\vec v,\vec w)$ forment une famille \textbf{liée} $\iff$ ils sont coplanaires.
\end{itemize}
\end{propriete}

\begin{exemple}[Étudier une famille de vecteurs]
Dans le plan muni d'un repère, soient $\vec u(2,3)$, $\vec v(-4,-6)$. Montrons qu'ils sont colinéaires.

\smallskip
On cherche $\lambda$ tel que $\vec v=\lambda\vec u$ : $-4=2\lambda$ et $-6=3\lambda$ donnent tous deux $\lambda=-2$. Donc $\vec v=-2\vec u$ : les vecteurs sont colinéaires (famille liée).

\smallskip
\emph{Critère par déterminant} : $\vec u(x,y)$, $\vec v(x',y')$ colinéaires $\iff xy'-x'y=0$. Ici : $2\times(-6)-(-4)\times3=-12+12=0$. $\checkmark$
\end{exemple}

% ------------------- SECTION 3 -------------------
\section{Sous-espaces vectoriels}

\begin{definition}[Sous-espace vectoriel]
Soit $E$ un espace vectoriel et $F\subset E$. $F$ est un \textbf{sous-espace vectoriel} de $E$ lorsque $F$, muni des lois restreintes, est lui-même un espace vectoriel.
\end{definition}

\begin{propriete}[Caractérisation pratique]
$F\subset E$, $F\ne\varnothing$, est un sous-espace vectoriel de $E$ si et seulement si :
\[
  (\forall \vec u,\vec v\in F)(\forall\lambda\in\mathbb{R})\quad \vec u+\lambda\vec v \in F.
\]
En particulier, un sous-espace vectoriel contient toujours $\vec0$.
\end{propriete}

\begin{demo}
Si $F$ est un sous-espace vectoriel, la stabilité par combinaison linéaire découle directement des axiomes (loi interne $+$ et loi externe). Réciproquement, si la propriété est vérifiée : en prenant $\vec v=\vec u$, $\lambda=-1$, on obtient $\vec u-\vec u=\vec0\in F$ ; les axiomes d'espace vectoriel se transmettent alors de $E$ à $F$ (associativité, commutativité, etc. sont automatiquement vérifiées puisque valables sur $E$ tout entier).
\end{demo}

\begin{exemple}[Droite vectorielle]
Soit $\vec u\ne\vec0$ un vecteur du plan. Montrons que $F=\{\lambda\vec u : \lambda\in\mathbb{R}\}$ (l'ensemble des vecteurs colinéaires à $\vec u$) est un sous-espace vectoriel.

\smallskip
$F\ne\varnothing$ ($\vec0=0\cdot\vec u\in F$). Soient $\vec x=\lambda\vec u$, $\vec y=\mu\vec u\in F$, et $\alpha\in\mathbb{R}$ :
\[
  \vec x+\alpha\vec y = \lambda\vec u+\alpha\mu\vec u = (\lambda+\alpha\mu)\vec u \in F.
\]
Donc $F$ est un sous-espace vectoriel de $\mathcal{V}$, appelé \textbf{droite vectorielle} dirigée par $\vec u$.
\end{exemple}

\begin{exemple}[Contre-exemple : ensemble non stable]
L'ensemble $G=\{(x,y)\in\mathbb{R}^2 : x^2+y^2=1\}$ (le cercle unité, vu comme ensemble de vecteurs) n'est \textbf{pas} un sous-espace vectoriel : par exemple $(1,0)\in G$ mais $2\times(1,0)=(2,0)\notin G$ (la loi externe ne stabilise pas $G$).
\end{exemple}

% ------------------- SECTION 4 -------------------
\section{Bases et coordonnées}

\begin{definition}[Base]
Une famille $(\vec e_1,\ldots,\vec e_n)$ est une \textbf{base} de l'espace vectoriel $E$ lorsqu'elle est \textbf{libre} et \textbf{génératrice} (tout vecteur de $E$ est combinaison linéaire des $\vec e_i$).
\end{definition}

\begin{propriete}[Bases du plan et de l'espace]
\begin{itemize}[leftmargin=6mm, itemsep=0.5mm]
  \item Deux vecteurs \textbf{non colinéaires} $(\vec i,\vec j)$ forment une \textbf{base du plan} : tout vecteur $\vec w$ du plan s'écrit de façon \textbf{unique} $\vec w=x\vec i+y\vec j$, $(x,y)$ étant les \textbf{coordonnées} de $\vec w$ dans cette base.
  \item Trois vecteurs \textbf{non coplanaires} $(\vec i,\vec j,\vec k)$ forment une \textbf{base de l'espace} : tout vecteur $\vec w$ s'écrit de façon unique $\vec w=x\vec i+y\vec j+z\vec k$.
\end{itemize}
\end{propriete}

\begin{propriete}[Unicité de la décomposition]
Dans une base, l'écriture d'un vecteur comme combinaison linéaire des vecteurs de base est \textbf{unique}.
\end{propriete}

\begin{demo}
Supposons $x\vec i+y\vec j=x'\vec i+y'\vec j$. Alors $(x-x')\vec i=(y'-y)\vec j$. Si $x\ne x'$, on obtiendrait $\vec i=\dfrac{y'-y}{x-x'}\vec j$, c'est-à-dire $\vec i$ et $\vec j$ colinéaires, ce qui contredit l'hypothèse que $(\vec i,\vec j)$ est une base (famille libre). Donc $x=x'$, puis $(y'-y)\vec j=\vec0$ avec $\vec j\ne\vec0$, donc $y=y'$.
\end{demo}

\begin{exemple}[Décomposition dans une base]
Dans une base $(\vec i,\vec j)$ du plan, soient $\vec u=2\vec i-\vec j$, $\vec v=\vec i+3\vec j$. Exprimons $\vec w=5\vec i+\vec j$ comme combinaison linéaire de $\vec u$ et $\vec v$.

\smallskip
On vérifie d'abord que $(\vec u,\vec v)$ est une base : $2\times3-1\times(-1)=7\ne0$, donc $\vec u,\vec v$ non colinéaires.

\smallskip
On cherche $a,b$ tels que $\vec w=a\vec u+b\vec v$ : $5\vec i+\vec j=a(2\vec i-\vec j)+b(\vec i+3\vec j)=(2a+b)\vec i+(-a+3b)\vec j$.

\smallskip
Par unicité de la décomposition dans la base $(\vec i,\vec j)$ : $2a+b=5$ et $-a+3b=1$. De la première : $b=5-2a$. En substituant : $-a+3(5-2a)=1\iff-7a+15=1\iff a=2$, d'où $b=1$.

\smallskip
\textbf{Conclusion} : $\vec w=2\vec u+\vec v$ (vérification : $2(2\vec i-\vec j)+(\vec i+3\vec j)=5\vec i+\vec j$ $\checkmark$).
\end{exemple}

% ------------------- SECTION 5 -------------------
\section{Exercices corrigés -- niveau examen national SM}

\begin{exemple}[Exercice 1 -- montrer qu'une famille est une base]
Dans un repère $(O,\vec i,\vec j)$, soient $\vec u(1,2)$ et $\vec v(3,1)$. Montrer que $(\vec u,\vec v)$ est une base du plan.

\smallskip
\textbf{Corrigé.} Il suffit de montrer que $\vec u$ et $\vec v$ ne sont pas colinéaires : $1\times1-3\times2=1-6=-5\ne0$. Donc $(\vec u,\vec v)$ est libre, et comme le plan est de dimension $2$, toute famille libre de $2$ vecteurs en est une base.
\end{exemple}

\begin{exemple}[Exercice 2 -- sous-espace vectoriel dans $\mathbb{R}^3$]
Soit $F=\{(x,y,z)\in\mathbb{R}^3 : x+y-z=0\}$. Montrer que $F$ est un sous-espace vectoriel de $\mathbb{R}^3$.

\smallskip
\textbf{Corrigé.} $F\ne\varnothing$ : $(0,0,0)\in F$ (car $0+0-0=0$).

\smallskip
Soient $(x,y,z),(x',y',z')\in F$ et $\lambda\in\mathbb{R}$. Montrons que $(x,y,z)+\lambda(x',y',z')=(x+\lambda x',y+\lambda y',z+\lambda z')\in F$ :
\[
  (x+\lambda x')+(y+\lambda y')-(z+\lambda z') = (x+y-z)+\lambda(x'+y'-z') = 0+\lambda\times0 = 0.
\]
Donc $F$ est un sous-espace vectoriel de $\mathbb{R}^3$ (un \textbf{plan vectoriel}, car il correspond aux vecteurs dont les coordonnées vérifient une équation linéaire homogène).
\end{exemple}

\begin{exemple}[Exercice 3 -- famille liée et coplanarité]
Dans l'espace, soient $\vec u(1,1,1)$, $\vec v(2,0,1)$, $\vec w(4,2,3)$. Montrer que ces trois vecteurs sont coplanaires en exprimant $\vec w$ en fonction de $\vec u$ et $\vec v$.

\smallskip
\textbf{Corrigé.} On cherche $a,b$ tels que $\vec w=a\vec u+b\vec v$ :
\[
  \begin{cases} a+2b=4 \\ a=2 \\ a+b=3 \end{cases}
\]
De la deuxième équation : $a=2$. En substituant dans la troisième : $b=1$. Vérification dans la première : $2+2\times1=4$. $\checkmark$

\smallskip
\textbf{Conclusion} : $\vec w=2\vec u+\vec v$ : les trois vecteurs sont coplanaires (famille liée).
\end{exemple}

% ------------------- SECTION 6 -------------------
\section{Méthodes et pièges : la boîte à outils de l'examen}

\subsection{Ce qu'on te demande, ce que tu fais}

\begin{center}
\renewcommand{\arraystretch}{1.45}
\sffamily\small
\begin{tabular}{>{\raggedright\arraybackslash}p{7.2cm} >{\raggedright\arraybackslash}p{8cm}}
\rowcolor{qtealdark} \color{white}\bfseries Ce qu'on te demande & \color{white}\bfseries Ton réflexe \\
\rowcolor{tealpale} Montrer que deux vecteurs sont colinéaires & Chercher $\lambda$ tel que $\vec v=\lambda\vec u$, ou calculer le déterminant $xy'-x'y=0$. \\
Montrer qu'une famille de 2 (ou 3) vecteurs est une base & Prouver qu'elle est libre (non colinéaires / non coplanaires). \\
\rowcolor{tealpale} Montrer que $F$ est un sous-espace vectoriel & $F\ne\varnothing$ (souvent $\vec0\in F$), puis $\vec u+\lambda\vec v\in F$ pour tous $\vec u,\vec v\in F$, $\lambda\in\mathbb{R}$. \\
Décomposer un vecteur dans une base & Poser le système linéaire issu de l'identification des coordonnées, résoudre. \\
\rowcolor{tealpale} Montrer que 3 vecteurs de l'espace sont coplanaires & Exprimer l'un comme combinaison linéaire des deux autres. \\
\end{tabular}
\end{center}

\subsection{Les pièges qui coûtent des points}

\begin{remarque}[Erreurs relevées chaque année par les correcteurs]
\begin{itemize}[leftmargin=6mm, itemsep=0.5mm]
  \item Oublier de vérifier $F\ne\varnothing$ avant d'utiliser la caractérisation pratique d'un sous-espace vectoriel.
  \item Confondre « famille libre » (aucune combinaison non triviale ne donne $\vec0$) et « famille génératrice » (elle engendre tout l'espace) : une base exige les \textbf{deux}.
  \item Oublier qu'un sous-espace vectoriel contient \textbf{toujours} $\vec0$ : un ensemble qui ne le contient pas n'en est jamais un.
  \item Confondre le critère de colinéarité $xy'-x'y=0$ avec une simple égalité de rapports mal posée quand une coordonnée est nulle.
  \item Oublier l'\textbf{unicité} de la décomposition dans une base pour justifier l'identification terme à terme des coordonnées.
\end{itemize}
\end{remarque}

% ------------------- RESUME -------------------
\vspace{4mm}
\begin{tcolorbox}[
  enhanced, boxrule=0pt, arc=2mm,
  interior style={left color=qtealdark, right color=qteal},
  left=6mm, right=6mm, top=4mm, bottom=4mm,
  fontupper=\color{white}\sffamily,
  title={\sffamily\bfseries\color{white}\ding{72}~~À retenir},
  colbacktitle=qtealdark, coltitle=white, boxrule=0pt
]
\begin{itemize}[leftmargin=6mm, label={\color{qamber}\ding{212}}, itemsep=1mm]
  \item L'ensemble des vecteurs du plan (resp. de l'espace), avec $+$ et $\cdot$, est un espace vectoriel réel.
  \item Famille libre $\iff$ aucune combinaison linéaire non triviale ne vaut $\vec0$ ; sinon, famille liée.
  \item Sous-espace vectoriel $F$ : $F\ne\varnothing$ et $\vec u+\lambda\vec v\in F$ pour tous $\vec u,\vec v\in F$, $\lambda\in\mathbb{R}$.
  \item Base : famille libre et génératrice ; décomposition d'un vecteur dans une base \textbf{unique}.
  \item Dans le plan : 2 vecteurs non colinéaires forment une base ; dans l'espace : 3 vecteurs non coplanaires.
\end{itemize}
\end{tcolorbox}

\end{document}
